\begin{solapartat}
\punts{0,25} L'equació d'una ona és: $\psi(x, t) = A\cos(\omega t - kx) = A\cos 2\pi\left(\dfrac{t}{T} - \dfrac{x}{\lambda}\right)$

% DUBTE: la pauta dona k en rad/s (hauria de ser rad/m).
\punts{0,25} Llavors $\omega = 2,40\times 10^{15}\ \mathrm{rad/s}$ i $k = 1,20\times 10^{7}\ \mathrm{rad/s}$
\begin{align*}
T &= \frac{2\pi}{\omega} = 2,618\times 10^{-15}\ \mathrm{s}\\
f &= \frac{1}{T} = \frac{\omega}{2\pi} = 3,82\times 10^{14}\ \mathrm{Hz}\\
\lambda &= \frac{2\pi}{k} = 5,24\times 10^{-7}\ \mathrm{m} \quad \text{\punts{0,1}}
\end{align*}

\punts{0,1} $v = \dfrac{\lambda}{T} = \lambda f = 2,00\times 10^{8}\ \mathrm{m/s}$

\punts{0,2} $n = \dfrac{c}{v} = 1,5$

\punts{0,15} En el vuit la velocitat de propagació és $c = 3,00\times 10^{8}\ \mathrm{m/s}$

% DUBTE: la pauta escriu λ = cf; el valor 7,85e-7 m correspon a λ = c/f.
\punts{0,2} $\lambda = cf = 7,85\times 10^{-7}\ \mathrm{m}$
\end{solapartat}

\begin{solapartat}
\punts{0,65} La direcció de propagació és l'eix $x$ atès que la coordenada $x$ apareix en l'argument de l'equació de l'ona.

\punts{0,6} El camp magnètic és perpendicular a la direcció de propagació (eix $x$) i al camp elèctric. Segons l'equació de l'ona electromagnètica el camp elèctric oscil·la en l'eix ($z$) (segons l'equació de l'ona, la direcció $\vec{E}(x, t)$ ve donada pel vector unitari $\vec{k}$). Per tant, el camp magnètic oscil·la en la direcció de l'eix $y$.
\end{solapartat}
